\documentclass[openany,fontset=fandol]{ctexbook}
\usepackage{amsmath,amssymb}
\usepackage[unicode,hidelinks]{hyperref}
\newcommand{\note}[1]{\footnote{译注：#1}}
\begin{document}
\title{代数几何基础第五卷预备笔记\\第12节：正则截面的推广说明}
\author{据用户上传英译编辑稿译出}
\date{2026年11月16日}
\maketitle
\chapter*{范围与校订状态}
本批连续译出底本第45页末四段，自“[Le laius 5.1 Fr or Latin 67] goes as such”起，止于“so does 5.7”。下一处是第46页“Proposition 6.1”。这是第12节中对第5节的推广指示，不是第5节各命题的完整重述或证明。本稿为独立续译，未合入主稿；此前缺文与待核事项不因本批而消除。
\chapter*{§12 将上述结果推广到线性截面（续）}
5.1中的那段论述可以照原样保留。\note{底本作“[Le laius 5.1 Fr or Latin 67]”，保留着英译编辑者对语言及页码的标记；这里按其语境译作“那段论述”，不把67当作数学条目号。}在5.2中，须记住：$\phi$现在成为$G_T(1)$的一个截面$\phi^{(m)}$，其中
\[
T=X\times_S\operatorname{Gr}^{m}.
\]
对于$\xi\in\operatorname{Gr}^{m}$，它诱导出
\[
\phi^{(m)}_{\xi}\in
\Gamma\bigl(X_{k(\xi)},\mathcal O_{X_{k(\xi)}}(1)\otimes_{k(\xi)}G(\xi)\bigr).
\]
\note{底本的$T$排作“$X_Sx\mathrm{Gr}^m$”，此处展开为纤维积；截面所在模排作$\mathcal O_{X_{k(\xi)}}(1)_{k(\xi)}\otimes G(\xi)$，重复的基变换下标在正文整理为通常写法。$G$仍指本节开头的泛秩$m$商模，$G(\xi)$是其在$\xi$处的纤维；并非第5节中作为限制模所用的同名字母。}

不过，一般地，我们将在第19段说明如下概念：给定一个预概形上的秩$m$局部自由模的截面$\phi$，以及一个模$F$，怎样称这一截面为$F$-正则。用局部基来表达，这意味着它的$m$个分量是$\mathcal O_X$的截面，并组成一个$F$-正则序列（还须验证这不依赖于所选的基）。\note{底本前半句写“rank $n$”，后半句却写“$m$ sections”；本段及相邻两段以$m$表示截面的秩。秩$m$的局部自由模在一组基下恰有$m$个分量，故正文将前一处$n$统一为$m$，而非改变所需序列的长度。底本还留有“this means that such a section if $F$ regular”的句法残损，按后文定义性语境译出。基无关性仍按原稿保留为待验证事项，不在这里宣称已给出证明。}当$m=1$时，就有5.2中提到的那个显然的内在解释。采用这一术语约定后，5.3仍可照原样保留，5.4也是如此。

\newpage
附注5.5的第一部分可以推广到任意$m$的情形：若$F_s$满足$(S_m)$，则关于$\phi_\xi$的上述正则性条件可以用纯粹的维数条件来表达。\note{底本此处排作$F_S$。所指5.5先固定$\xi$在$s\in S$之上，再假设纤维模$F_s$满足$(S_1)$；此处推广的是同一个纤维上的条件，故正文写$F_s$。原稿没有在本段给出具体维数判据，本批不凭概括性指示补造公式或删减前文假设。}

附注5.5的第二部分对于任意$m$都可照原样保留。定理5.6也照原样推广，5.7也是如此。\note{“照原样”是原作者的推广指示，并不意味着可以丢弃5.6中的有限表示、相对于$S$平坦等假设或5.7中的固有性条件；本批仅译出指示本身，未重写这些条目的完整陈述和证明。}
\end{document}
